% chapters/implementation/methodology.tex
% Sección 5.1: Metodología de desarrollo (~1.5-2 páginas)
% ============================================================================
%
% GUÍA: Describir cómo se organizó el desarrollo del proyecto.
%
%   1. Metodología elegida: desarrollo iterativo e incremental.
%      - Reuniones semanales con el tutor (Eloy Anguiano Rey) y con el
%        compañero del TFG de frontend (Alejandro Soler Sichling).
%      - Cada iteración produce un incremento funcional del backend.
%      - Elementos de Scrum adaptados al contexto de un TFG:
%        * Planificación semanal (qué subsistema implementar).
%        * Revisión con el tutor (demostración del incremento).
%        * No hay retrospectiva formal (equipo de 2 personas).
%
%   2. Fases del proyecto (descripción textual, NO Gantt):
%      Describir las fases/iteraciones principales:
%        a) Fase de análisis y diseño inicial:
%           - Definición de requisitos.
%           - Diseño de la arquitectura y modelo de datos.
%           - Configuración del entorno de desarrollo (Docker Compose).
%        b) Fase de implementación del núcleo:
%           - Modelos de datos y migraciones.
%           - Sistema de autenticación pluggable.
%           - Sistema de permisos.
%           - CRUD de entidades académicas.
%        c) Fase de implementación de subsistemas avanzados:
%           - Pipeline OCR e ingesta de PDFs.
%           - Corrección y anotaciones multimedia.
%           - Máquina de estados del examen.
%           - Notificaciones y correo.
%        d) Fase de pruebas y refinamiento:
%           - Pruebas unitarias y de integración.
%           - Pruebas de carga.
%           - Correcciones y optimizaciones.
%        e) Fase de documentación:
%           - Redacción de la memoria.
%
%   --- DIAGRAMA DE GANTT (PLACEHOLDER DESCARTABLE) ---
%   Si finalmente decides incluir un Gantt, descomenta lo siguiente y
%   ajusta las semanas. Si no, elimina este bloque.
%
%   % \begin{figure}{fig:gantt}{Planificación temporal del proyecto}
%   %   \begin{gantt}{1}{20}
%   %     \taskgroup[g1]{Análisis y diseño}{1}{4}
%   %     \taskbar[t1]{Requisitos}{1}{2}
%   %     \taskbar[t2]{Diseño de arquitectura}{2}{4}
%   %     \\
%   %     \taskgroup[g2]{Implementación núcleo}{4}{10}
%   %     \taskbar[t3]{Modelos y auth}{4}{6}
%   %     \taskbar[t4]{Permisos y CRUD}{6}{8}
%   %     \taskbar[t5]{API REST base}{8}{10}
%   %     \\
%   %     \taskgroup[g3]{Subsistemas avanzados}{10}{16}
%   %     \taskbar[t6]{Pipeline OCR}{10}{13}
%   %     \taskbar[t7]{Corrección}{12}{15}
%   %     \taskbar[t8]{Ciclo de vida}{14}{16}
%   %     \\
%   %     \taskgroup[g4]{Pruebas y documentación}{15}{20}
%   %     \taskbar[t9]{Pruebas}{15}{18}
%   %     \taskbar[t10]{Memoria}{16}{20}
%   %     \\
%   %     \milestone[m1]{Entrega}{20}
%   %     \FtoSlink{t2}{t3}
%   %     \FtoSlink{t5}{t6}
%   %     \FtoSlink{t8}{t9}
%   %   \end{gantt}
%   % \end{figure}
%
%   3. Herramientas de desarrollo:
%      - Control de versiones: Git (GitHub/GitLab).
%      - IDE: VS Code con extensiones de Python/Django.
%      - Entorno: Docker Compose para desarrollo local.
%      - TODO: Añadir CI/CD si lo configuraste.
%
%   4. Coordinación con el TFG de frontend:
%      - Contrato de API compartido (especificación OpenAPI).
%      - Reuniones conjuntas para definir endpoints.
%      - Desarrollo paralelo con integración periódica.
% ============================================================================

% TODO: Redactar la sección de metodología.
